\documentclass{article}
\usepackage{amsmath}
\usepackage{xcolor}
\begin{document}

SERIES DE POTENCIA

Son una representacion a traves de una combinacion de teminos, que comunemente siguen un patron en general como:




\begin{gather*}
1+x+x^{2}+x^{3}+.....x^{n}   \rightarrow \text{donde n, } \\
\text{representa la posicon o el indice}
\end{gather*}


Donde en dicha representacion, dado el valor de x, se puede determinar un valor como tal.

Existen otros tipos de series matematicas con las series de fibonnaci:




\begin{gather*}
1 1 2  3  5  8  13  ...... \\
\text{donde se sabe que f}_{n} = f_{n-1}+f_{n-2}
\end{gather*}


En el contexto de E.D. una combinacion de terminos representados por una serie pueden ser una posible solucion a un E.D.

En una E.D. por decir, se puede asumir que una solucion podria ser : 
\begin{gather*}
y=c_{1}e^{x}
\end{gather*}
, pero con series de potencia, se podria asumir una representacion alternativa como:




\begin{gather*}
y =c_{1}(1+x + \frac{x^{2}}{2!}+\frac{x^{3}}{3!}+\frac{x^{4}}{4!}+....)
\end{gather*}


supongamos que x = 0.7 y c1 = 1

en la supocion donde 
\begin{gather*}
y=c_{1}e^{x}= 1e^{0.7}=2.01 
\end{gather*}


Si asumimos la representacion por series de potencia... deberiamos llegar un resultado aproximado

Consideremos tomar inicialmente 2 terminos de la serie: 


\begin{gather*}
y =c_{1}(1+x) \\
y = 1(1+0.7)= 1.7
\end{gather*}


Consideremos tomar adicionalmente 2 terminos mas:


\begin{gather*}
y =c_{1}(1+x + \frac{x^{2}}{2!}+\frac{x^{3}}{3!}) \\
y = 1(1+0.7+\frac{0.7^{2}}{2}+\frac{0.7^{3}}{6})=2.002 \\
             0.24 + 0.05
\end{gather*}


Consideremos tomar adicionalmente 2 terminos mas:


\begin{gather*}
y =c_{1}(1+x + \frac{x^{2}}{2!}+\frac{x^{3}}{3!}+\frac{x^{4}}{4!}+\frac{x^{5}}{5!}) \\
y = 1(1+0.7+\frac{0.7^{2}}{2}+\frac{0.7^{3}}{6}+\frac{0.7^{4}}{24}+\frac{0.7^{5}}{120})=2.01
\end{gather*}


Las series de potencia son util para cuando los patrones de E.D. vistos en anteriores capitulos no encajan en los ejercicios dados, por ejemplo


\begin{gather*}
y''+5y'-8\text{\colorbox[RGB]{248,231,28}{$x$}}y=0   \rightarrow \text{ Simple vista, pareceria una E.D. con CC H.} \\
\text{pero la x resaltada, no permite aplicar dicho metodo} \\
\vspace{0.5em} \\
x^{3}y'''+5x^{2}y''-5\text{\colorbox[RGB]{248,231,28}{$x$}}\text{\colorbox[RGB]{248,231,28}{$^{2}$}}y+2y=0  \rightarrow \text{el termino resltado ya no permite} \\
\text{que la E.D. sea resuelta por Cauchy Euler}
\end{gather*}


Para resolver algunos casos, es posible recurrir a utilizar series de Potencia, donde se asume que la solucion a la E.D. esta representada por:


\begin{gather*}
y = a_{0}+a_{1}(x-x_{0})+a_{2}(x^{2}-x_{0})+a_{3}(x^{3}-x_{0})+... \\
\text{donde a}_{0},a_{1},a_{2}\text{,etc son coeficientes}
\end{gather*}


Es representacion esta centrada en 
\begin{gather*}
x_{0}
\end{gather*}
, para simplificar la representacion, podemos asumir que 
\begin{gather*}
x_{0}=0
\end{gather*}
, en ese sentido se tiene:


\begin{gather*}
y = a_{0}+a_{1}x+a_{2}x^{2}+a_{3}x^{3}+ a_{4}x^{4}...
\end{gather*}


Si nuestra ecuacion diferencial es de 2do orden$\rightarrow$  sabemos que deberian existir 2 soluciones, acompañadas por una c1, y una c2. En ese sentido la forma de representar la solucion a la E.D. parte de \colorbox[RGB]{248,231,28}{poner en funcion una ctte de indice superior en funcion a una de indice inferior}.


\begin{gather*}
y = \text{\colorbox[RGB]{248,231,28}{$a$}}\text{\colorbox[RGB]{248,231,28}{$_{0}$}}+\text{\colorbox[RGB]{80,227,194}{$a$}}\text{\colorbox[RGB]{80,227,194}{$_{1}$}}\text{\colorbox[RGB]{80,227,194}{$x$}}+\text{\colorbox[RGB]{248,231,28}{$($}}\text{\colorbox[RGB]{248,231,28}{$ $}}\text{\colorbox[RGB]{248,231,28}{$\frac{1}{2}$}}\text{\colorbox[RGB]{248,231,28}{$a$}}\text{\colorbox[RGB]{248,231,28}{$_{0}$}}\text{\colorbox[RGB]{248,231,28}{$)$}}\text{\colorbox[RGB]{248,231,28}{$x$}}\text{\colorbox[RGB]{248,231,28}{$^{2}$}}+\text{\colorbox[RGB]{80,227,194}{$3a$}}\text{\colorbox[RGB]{80,227,194}{$_{1}$}}\text{\colorbox[RGB]{80,227,194}{$x$}}\text{\colorbox[RGB]{80,227,194}{$^{3}$}}+ \\
\text{Supongamos que es posible deducir que a}_{2}, y a_{3}\text{ se pueden representar en funcion a a}_{0} y a_{1} \\
a_{2}= \frac{1}{2}a_{0} \\
a3 = 3a_{1}
\end{gather*}





\begin{gather*}
y''+5y'+6y=0 \\
(r+2)(r+3)=0 \\
y = c_{1}e^{-2x}+c_{2}e^{-3x} = c_{1}y_{1}+c_{2}y_{2}
\end{gather*}


Ahora, en toda E.D. para demostrar que una solucion esta acorde a la E.D.,


\begin{gather*}
y = a_{0}+a_{1}x+a_{2}x^{2}+a_{3}x^{3}+ a_{4}x^{4}+... \\
y' = a_{1}+2a_{2}x+3a_{3}x^{2}+ 4a_{4}x^{3} \\
y'' = 2a_{2}+6a_{3}x+ 12a_{4}x^{2}
\end{gather*}


Supongamos que nos facilitan la siguiente E.D.


\begin{gather*}
y''+3xy'-5y=0 \\
2a_{2}+6a_{3}x+ 12a_{4}x^{2}+3x(a_{1}+2a_{2}x+3a_{3}x^{2}+ 4a_{4}x^{3}) \\
-5(a_{0}+a_{1}x+a_{2}x^{2}+a_{3}x^{3}+ a_{4}x^{4})=0 \\
\vspace{0.5em} \\
\text{\colorbox[RGB]{248,231,28}{$2a$}}\text{\colorbox[RGB]{248,231,28}{$_{2}$}}+\text{\colorbox[RGB]{245,166,35}{$6a$}}\text{\colorbox[RGB]{245,166,35}{$_{3}$}}\text{\colorbox[RGB]{245,166,35}{$x$}}+ 12a_{4}x^{2}+\text{\colorbox[RGB]{245,166,35}{$3a$}}\text{\colorbox[RGB]{245,166,35}{$_{1}$}}\text{\colorbox[RGB]{245,166,35}{$x$}}+6a_{2}x^{2}+9a_{3}x^{3}+ 12a_{4}x^{4} \\
\text{\colorbox[RGB]{248,231,28}{$-5a$}}\text{\colorbox[RGB]{248,231,28}{$_{0}$}}\text{\colorbox[RGB]{245,166,35}{$-5a$}}\text{\colorbox[RGB]{245,166,35}{$_{1}$}}\text{\colorbox[RGB]{245,166,35}{$x$}}-5a_{2}x^{2}-5a_{3}x^{3}-5 a_{4}x^{4}=0 \\
(2a_{2}-5a_{0}) = 0   \rightarrow a_{2} =\frac{5}{2}a_{0} \\
(6a_{3}x+3a_{1}x-5a_{1}x)=0  \rightarrow x(6a_{3}-2a_{1})=0  \rightarrow a_{3} =\frac{1}{3}a_{1} \\
\vspace{0.5em} \\
(12a_{4}x^{2}+6a_{2}x^{2}-5a_{2}x^{2})=0  \rightarrow x^{2}(12a_{4}+a_{2}) \implies  a_{4}=-\frac{1}{12}a_{2}=-\frac{1}{12}*\frac{5}{2}a_{0} \\
(9a_{3}x^{3}-5a_{3}x^{3}) =0  \rightarrow \text{no me sirve.. porque no se tomo mas terminos} \\
(12a_{4}x^{4}-5 a_{4}x^{4})=0 \rightarrow \text{no me sirve.. porque no se tomo mas terminos} \\
y = a_{0}+a_{1}x+a_{2}x^{2}+a_{3}x^{3}+ a_{4}x^{4}+... \\
y =\text{\colorbox[RGB]{248,231,28}{$ a$}}\text{\colorbox[RGB]{248,231,28}{$_{0}$}}+a_{1}x+\text{\colorbox[RGB]{248,231,28}{$\frac{5}{2}$}}\text{\colorbox[RGB]{248,231,28}{$a$}}\text{\colorbox[RGB]{248,231,28}{$_{0}$}}\text{\colorbox[RGB]{248,231,28}{$x$}}\text{\colorbox[RGB]{248,231,28}{$^{2}$}}+\frac{1}{3}a_{1}x^{3}\text{\colorbox[RGB]{248,231,28}{$- $}}\text{\colorbox[RGB]{248,231,28}{$\frac{5}{24}$}}\text{\colorbox[RGB]{248,231,28}{$a$}}\text{\colorbox[RGB]{248,231,28}{$_{0}$}}\text{\colorbox[RGB]{248,231,28}{$x$}}\text{\colorbox[RGB]{248,231,28}{$^{4}$}}+... \\
y=a_{0}(1+\frac{5}{2}x^{2}-\frac{5}{24}x^{4}+...)+a_{1}(x+\frac{1}{3}x^{3}+...) \\
\vspace{0.5em} \\
y_{1} =1+\frac{5}{2}x^{2}-\frac{5}{24}x^{4}+... \\
y_{2} = x+\frac{1}{3}x^{3}+... \\
\text{Intentemos demostrar que son aproximadamente una solucion a la E.D} \\
y''+3xy'-5y=0 \\
\vspace{0.5em} \\
y_{1} =1+\frac{5}{2}x^{2}-\frac{5}{24}x^{4}+... \\
y_{1}' =5x-\frac{5}{6}x^{3} \\
y_{1}'' =5-\frac{5}{2}x^{2} \\
\vspace{0.5em} \\
5-\frac{5}{2}x^{2}+3x(5x-\frac{5}{6}x^{3})-5(1+\frac{5}{2}x^{2}-\frac{5}{24}x^{4})=0 \\
5\text{\colorbox[RGB]{248,231,28}{$-$}}\text{\colorbox[RGB]{248,231,28}{$\frac{5}{2}$}}\text{\colorbox[RGB]{248,231,28}{$x$}}\text{\colorbox[RGB]{248,231,28}{$^{2}$}}\text{\colorbox[RGB]{248,231,28}{$+15x$}}\text{\colorbox[RGB]{248,231,28}{$^{2}$}}\text{\colorbox[RGB]{126,211,33}{$-$}}\text{\colorbox[RGB]{126,211,33}{$\frac{5}{2}$}}\text{\colorbox[RGB]{126,211,33}{$x$}}\text{\colorbox[RGB]{126,211,33}{$^{4}$}}-5\text{\colorbox[RGB]{248,231,28}{$-$}}\text{\colorbox[RGB]{248,231,28}{$\frac{25}{2}$}}\text{\colorbox[RGB]{248,231,28}{$x$}}\text{\colorbox[RGB]{248,231,28}{$^{2}$}}\text{\colorbox[RGB]{126,211,33}{$+$}}\text{\colorbox[RGB]{126,211,33}{$\frac{25}{24}$}}\text{\colorbox[RGB]{126,211,33}{$x$}}\text{\colorbox[RGB]{126,211,33}{$^{4}$}}=0 \\
\vspace{0.5em} \\
\text{Efectivamente los terminos marcados con verde, no satisfacen al 100% la igualdad, ya } \\
-1.45  \rightarrow \text{ no es igual que cero, sin embargo, si se asumirian mas terminos} \\
el -1.45  se reduciria, y podria aproximarse a ser una solucion valida. \\
\vspace{0.5em} \\
\vspace{0.5em} \\
\text{Recordar Coeficientes Indeterminados para un polinomio} \\
y''+5y'+6y=x^{3}+2 \\
y_{p} = Ax^{3}+Bx^{2}+Cx+D \\
yp' =3Ax^{2}+2Bx^{}+C  \\
yp'' =6Ax+2B \\
6Ax+2B+15Ax^{2}+10Bx^{}+5C +Ax^{3}+Bx^{2}+Cx+D= x^{3}+2 \\
Ax^{3} = x^{3}   \rightarrow A=1 \\
15Ax^{2}+Bx^{2} =0 \\
6Ax+10Bx+Cx=0
\end{gather*}


Resolucion del ejercicio utilizando series generalizadas


\begin{gather*}
y''+3xy'-5y=x \\
\text{Se asume que una funcion puede estar representada de la siguiente forma} \\
y = \sum{∞}{n=0}a_{n}x^{n}=a_{0}+a_{1}x+a_{2}x^{2}+a_{3}x^{3}+ a_{4}x^{4}+...∞ \\
y' = \sum{∞}{n=1}na_{n}x^{n-1} =a_{1}+2a_{2}x+3a_{3}x^{2}+ 4a_{4}x^{3}+... \\
y' = \sum{∞}{n=2}n(n-1)a_{n}x^{n-2}= 2a_{2}+6a_{3}x+ 12a_{4}x^{2}+... \\
\sum{∞}{n=2}n(n-1)a_{n}x^{n-2}+\text{\colorbox[RGB]{248,231,28}{$3x$}}\sum{∞}{n=1}na_{n}x^{n-1}-5\sum{∞}{n=0}a_{n}x^{n}=x \\
\text{incluir cualquier valor de x}^{a}\text{ dentro de la seria, donde a es un real} \\
\sum{∞}{n=2}n(n-1)a_{n}x^{n-2}+\text{\colorbox[RGB]{248,231,28}{$3$}}\sum{∞}{n=1}na_{n}x^{n}-5\sum{∞}{n=0}a_{n}x^{n}=x \\
x^{0}                    x^{1}                   x^{0} \\
\text{Unificar los exponentes} \\
2a_{2}+ \sum{∞}{n=3}n(n-1)a_{n}x^{n-2}+\text{\colorbox[RGB]{248,231,28}{$3$}}\sum{∞}{n=1}na_{n}x^{n}-5a_{0}-5\sum{∞}{n=1}a_{n}x^{n}=x \\
\text{Reordenamos los elementos sueltos} \\
2a_{2}-5a_{0}+\sum{∞}{n=3}n(n-1)a_{n}x^{n-2}\text{\colorbox[RGB]{248,231,28}{$3$}}\sum{∞}{n=1}na_{n}x^{n}-5\sum{∞}{n=1}a_{n}x^{n}=x \\
2a_{2}-5a_{0}= 0  \rightarrow a2=\frac{5}{2}a_{0} \\
\sum{∞}{n=3}n(n-1)a_{n}x^{n-2}\text{\colorbox[RGB]{248,231,28}{$+3$}}\sum{∞}{n=1}na_{n}x^{n}-5\sum{∞}{n=1}a_{n}x^{n} \\
\text{Deseo que los tres terminos de series esten elevados a la 'k'} \\
\sum{∞}{n=3}n(n-1)a_{n}x^{n-2} = \sum{∞}{k+2=3}(k+2)(k+2-1)a_{k+2}x^{k} =\sum{∞}{k=1}(k+2)(k+1)a_{k+2}x^{k}  \\
k =n-2 \\
n=k+2 \\
\vspace{0.5em} \\
\sum{∞}{n=3}n(n-1)a_{n}x^{n-2}\text{\colorbox[RGB]{248,231,28}{$+3$}}\sum{∞}{n=1}na_{n}x^{n}-5\sum{∞}{n=1}a_{n}x^{n} \\
\vspace{0.5em} \\
k =n-2         k=n           k =n \\
n=k+2                                         \\
\sum{∞}{k=1}(k+2)(k+1)a_{k+2}x^{k} +3\sum{∞}{k=1}ka_{k}x^{k}-5\sum{∞}{k=1}a_{k}x^{k} \\
\vspace{0.5em} \\
\text{Como todos los terminos arrancan en un mismo indices, y tiene un mismo valor de x}^{k}, \\
\text{se puede factorizar x}^{k} \\
x^{k}[\sum{∞}{k=1}((k+2)(k+1)a_{k+2}+3ka_{k}-5a_{k})]=x \\
si k=1 \\
x[6a_{3}+3a_{1}-5a_{1} ] =x \\
6a_{3}-2a_{1}=1 \\
a_{3}=\frac{1+2a_{1}}{6}=\frac{1}{6}+\frac{1}{3}a_{1} \\
\vspace{0.5em} \\
\text{El resto de indices se puede directamente asumir que tendran igualdad 0} \\
para k>1 \\
(k+2)(k+1)a_{k+2}+(3k-5)a_{k}=0 \\
a_{k+2} =-\frac{(3k-5)}{(k+2)(k+1)}a_{k} \\
si k=2 \\
a_{4} = -\frac{(3*2-5)}{(2+2)(2+1)}a_{2}=-\frac{1}{12}a_{2} = -\frac{1}{12}*\frac{5}{2}a_{0} =-\frac{5}{24}a_{0} \\
\vspace{0.5em} \\
si k = 3 \\
a_{5}=-\frac{(3*3-5)}{(3+2)(3+1)}a_{3} =-\frac{4}{20}a_{3}=-\frac{4}{20}(\frac{1+2a_{1}}{6})=-\frac{1+2a_{1}}{30}=-\frac{1}{30}-\frac{a_{1}}{15} \\
\vspace{0.5em} \\
si k= 4 \\
a_{6}=-\frac{(3*4-5)}{(4+2)(4+1)}a_{4} =-\frac{7}{30}a_{4}=-\frac{7}{30}(-\frac{1}{12}*\frac{5}{2}a_{0})=\frac{7}{144}a_{0} \\
\vspace{0.5em} \\
si k=5 \\
a_{7}=-\frac{(3*5-5)}{(5+2)(5+1)}a_{5} =-\frac{10}{42}a_{5}=-\frac{10}{42}(-\frac{1}{30}-\frac{a_{1}}{15})=\frac{1}{126}+\frac{1a_{1}}{63} \\
\vspace{0.5em} \\
\text{Partimos de la solucion planteada anteriormente} \\
y = \sum{∞}{n=0}a_{n}x^{n}=a_{0}+a_{1}x+a_{2}x^{2}+a_{3}x^{3}+ a_{4}x^{4}+...∞ \\
\vspace{0.5em} \\
y =a_{0}+a_{1}x+\frac{5}{2}a_{0}x^{2}+(\frac{1}{6}+\frac{1}{3}a_{1})x^{3} -\frac{5}{24}a_{0}x^{4} \\
+(-\frac{1}{30}-\frac{a_{1}}{15})x^{5}+\frac{7}{144}a_{0}x^{6}+(\frac{1}{126}+\frac{a_{1}}{63})x^{7}+...∞ \\
\vspace{0.5em} \\
y''+3xy'-5y=x \\
y = y_{h}+y_{p} \\
y =\text{\colorbox[RGB]{248,231,28}{$a$}}\text{\colorbox[RGB]{248,231,28}{$_{0}$}}+\text{\colorbox[RGB]{126,211,33}{$a$}}\text{\colorbox[RGB]{126,211,33}{$_{1}$}}\text{\colorbox[RGB]{126,211,33}{$x$}}+\text{\colorbox[RGB]{248,231,28}{$\frac{5}{2}$}}\text{\colorbox[RGB]{248,231,28}{$a$}}\text{\colorbox[RGB]{248,231,28}{$_{0}$}}\text{\colorbox[RGB]{248,231,28}{$x$}}\text{\colorbox[RGB]{248,231,28}{$^{2}$}}+\frac{1}{6}x^{3}+\text{\colorbox[RGB]{126,211,33}{$\frac{1}{3}$}}\text{\colorbox[RGB]{126,211,33}{$a$}}\text{\colorbox[RGB]{126,211,33}{$_{1}$}}\text{\colorbox[RGB]{126,211,33}{$x$}}\text{\colorbox[RGB]{126,211,33}{$^{3}$}} \text{\colorbox[RGB]{248,231,28}{$-$}}\text{\colorbox[RGB]{248,231,28}{$\frac{5}{24}$}}\text{\colorbox[RGB]{248,231,28}{$a$}}\text{\colorbox[RGB]{248,231,28}{$_{0}$}}\text{\colorbox[RGB]{248,231,28}{$x$}}\text{\colorbox[RGB]{248,231,28}{$^{4}$}} \\
-\frac{1}{30}x^{5}\text{\colorbox[RGB]{126,211,33}{$-$}}\text{\colorbox[RGB]{126,211,33}{$\frac{1}{15}$}}\text{\colorbox[RGB]{126,211,33}{$a$}}\text{\colorbox[RGB]{126,211,33}{$_{1}$}}\text{\colorbox[RGB]{126,211,33}{$x$}}\text{\colorbox[RGB]{126,211,33}{$^{5}$}}+\text{\colorbox[RGB]{248,231,28}{$\frac{7}{144}$}}\text{\colorbox[RGB]{248,231,28}{$a$}}\text{\colorbox[RGB]{248,231,28}{$_{0}$}}\text{\colorbox[RGB]{248,231,28}{$x$}}\text{\colorbox[RGB]{248,231,28}{$^{6}$}}+\frac{1}{126}x^{7}\text{\colorbox[RGB]{126,211,33}{$+$}}\text{\colorbox[RGB]{126,211,33}{$\frac{1}{126}$}}\text{\colorbox[RGB]{126,211,33}{$a$}}\text{\colorbox[RGB]{126,211,33}{$_{1}$}}\text{\colorbox[RGB]{126,211,33}{$x$}}\text{\colorbox[RGB]{126,211,33}{$^{7}$}}
\end{gather*}





\end{document}
